\documentclass[11pt]{article}
\usepackage{amsmath,amssymb,amsthm,mathtools,mathrsfs}
\usepackage[margin=1in]{geometry}
\usepackage{enumitem}
\usepackage{hyperref}

\newtheorem{theorem}{Theorem}[section]
\newtheorem{proposition}[theorem]{Proposition}
\newtheorem{lemma}[theorem]{Lemma}
\newtheorem{corollary}[theorem]{Corollary}
\theoremstyle{definition}
\newtheorem{definition}[theorem]{Definition}
\newtheorem{construction}[theorem]{Construction}
\newtheorem{remark}[theorem]{Remark}
\newtheorem{example}[theorem]{Example}

\newcommand{\gp}{\mathrm{gp}}
\newcommand{\op}{\mathrm{op}}
\newcommand{\Sh}{\operatorname{Sh}}
\newcommand{\Lan}{\operatorname{Lan}}
\newcommand{\sat}{\mathrm{sat}}
\newcommand{\rk}{\operatorname{rk}}
\newcommand{\colim}{\operatorname*{colim}}
\newcommand{\HH}{\check{H}}
\newcommand{\cC}{\check{C}}
\newcommand{\Cone}{\mathrm{Cone}}
\newcommand{\XiLe}{\le^\Xi}
\newcommand{\XiLt}{<^\Xi}
\newcommand{\N}{\mathbb{N}}
\newcommand{\Z}{\mathbb{Z}}
\newcommand{\Q}{\mathbb{Q}}
\newcommand{\R}{\mathbb{R}}

\title{Problem 1 -- Humanize}
\author{}
\date{}

\begin{document}
\maketitle

\section{Statement}\label{sec:intro}

A \emph{monoid} is an integral, saturated, sharp, commutative, unital
semigroup, written additively, exactly as in the problem statement; it
induces the translation invariant partial order $x\ge0\iff x\in M$ on
$M^{\gp}$ and is \emph{valuative} if that order is total.
\emph{Infinitesimal} maps, \emph{birhizomicity at} a map, and
\emph{birhizomic covers} are used exactly as defined in the problem
statement. For a monoid $P$: $A_P$ is the category of monoid
homomorphisms $P\to Q$; $B_P\subseteq A_P$ the full subcategory of the
objects $P\to P[-c]^{\sharp}$; $A_P^{\op}$ carries the birhizomic
topology, $B_P^{\op}$ the trivial topology; the inclusion
$i:B_P^{\op}\subset A_P^{\op}$ induces
$\pi^{-1}:\Sh(B_P^{\op})\to\Sh(A_P^{\op})$ with right adjoint $\pi_{*}$.

\noindent\textbf{Scope hypothesis.} We prove the theorem for
\emph{finitely generated} $P$; from \S2 on, ``monoid'' means
finitely generated monoid wherever the finiteness of the face poset,
Gordan's lemma, or pointed-cone separation is used (Proposition~2.1,
Lemma~2.4, Corollary~2.5, and the face induction of \S4). Note that
every piece of every birhizomic cover of any object of $A_P$ is
\emph{a priori} finitely presented over $P$ (it is part of the
definition of birhizomicity at a valuative target), hence finitely
generated when $P$ is; the f.g.\ hypothesis is used precisely at the
places named above, and the general (non-finitely-generated) case is
obstructed exactly there: $P\to P/K$ need not be finitely presented
for arbitrary faces $K$, and the face poset can be infinite (see
Remark~\ref{rem:nonfg}).

We prove the degree-zero half of the desired quasi-isomorphism for
finitely generated $P$, and we record precisely why the higher-degree
half (and the general, non-finitely-generated case) remain open; see
the Closing remark (after Proposition~\ref{thm:main-end}).

\begin{theorem}\label{thm:main}
Let $P$ be a \emph{finitely generated} monoid \textup{(}see the Scope
hypothesis above\textup{)} and $F$ a presheaf of abelian groups on
$B_P^{\op}$. Then $\pi^{-1}F=G_F$ is a sheaf, and the degree-zero
unit
\[
F\;\xrightarrow{\ \sim\ }\;\pi_{*}\pi^{-1}F
\;=\;\mathrm R^{0}\pi_{*}\pi^{-1}F
\]
is an isomorphism. \textup{(}The higher direct images
$\mathrm R^{q>0}\pi_{*}(\pi^{-1}F)$ are \emph{not} shown to vanish
for arbitrary $F$; see the Closing remark.\textup{)}
\end{theorem}

\noindent\textbf{Overview.}
The argument rests on a collapse of the birhizome to dimension one
(\S\ref{sec:collapse}): every valuative monoid is $\{0\}$ or $\mathbb N$
(Lemma~\ref{lem:N}); hence valuative points are functionals
$P\to\mathbb N$ with half-space cones; and the birhizomic lifting
condition at a nonzero target collapses to one test pair
(Lemma~\ref{lem:atomic}, \emph{atomicity}), giving the \emph{normal
form} (Theorem~\ref{thm:normalform}): every cover refines to one by
\emph{folds} $(P\to\mathbb N)$ and \emph{CFA pieces} $(P\to P/K)$.
\S\ref{sec:sheaf} identifies $\pi^{-1}F$ with the explicit presheaf
$G_F=F\circ B$ (Proposition~\ref{prop:adj}) and proves it is a sheaf:
separatedness is formal; the descent computation
(Theorem~\ref{thm:descent}) boils down to the fact that the overlap of
any tight fold with the base piece $i(B(T))$ has ungrown kernel
(Proposition~\ref{prop:equiface}, Example~\ref{ex:ff}), so values are
transported by identity transitions. \S\ref{sec:direct} analyses the
higher sheaf cohomology on $i(b)$: the probe covers
(Definition~\ref{def:square}) are cofinal, the kernel computations of
the probe nerve (Lemma~\ref{lem:nerve}) split the \v{C}ech complex
into a fold strand and face strands, the fold terraces are killed
(Lemma~\ref{lem:sieve}) and the dead sector is made coboundary-valued;
the remaining residual computation is carried out as far as the
structure of a general presheaf permits, and its limitation for
arbitrary $F$ is recorded honestly in the Closing remark.

\paragraph{Notation.} For $T=(P\xrightarrow{\alpha_T}Q_T)\in A_P$:
$\Phi_T:=\ker(\alpha_T)$ (a face, Proposition~\ref{prop:bp}(b));
$B(T):=(P\to P/\Phi_T)\in B_P$;
\[
Z_T:=\{x\in P^{\gp}:\alpha_T^{\gp}(x)\in Q_T\},\qquad
L_T:=Z_T\cap(-Z_T),
\]
the \emph{relative saturation} and its \emph{lineality};
$P\cap L_T=\Phi_T$ (Lemma~\ref{lem:Zpointed}). For a face $K$: $P/K$
the face quotient (Proposition~\ref{prop:bp}(a));
$K^{\gp}\subseteq P^{\gp}$ its generated group; $K^{\gp}\cap P=K$.
A \emph{functional} is $\varphi:M\to\mathbb N$; $[\varphi]$ its
homothety class; ${GF}(\varphi)$ the positive generator of
$\varphi^{\gp}(M^{\gp})\subseteq\Z$; $\varphi$ is \emph{tight} on a
sharp $M$ if $\ker\varphi=\{0\}$. $\cC^\bullet,\HH^\bullet$ are
alternating \v{C}ech complexes/cohomology.

\section{Faces and the valuative collapse}\label{sec:collapse}

\subsection{Faces; the category $B_P$}

\begin{proposition}\label{prop:bp}
Let $M$ be a monoid, $K\subseteq M$ a face.
\begin{enumerate}[label=\textup{(\alph*)},leftmargin=1.8em]
\item $M/K:=(M\setminus K)\cup\{0\}$, with $0$ absorbing, is a monoid;
$M\to M/K$ is initial among maps killing $K$; and
$M/K\cong M[-c]^{\sharp}$ whenever $c$ generates the face $K$; in
particular $P[-c]^{\sharp}=P/\langle c\rangle_{\mathrm{face}}$, so
$B_P$ is the poset of faces of $P$: a morphism $P/K_1\to P/K_2$
\textup{(}unique\textup{)} exists iff $K_1\subseteq K_2$.
\item Faces are exactly the sets $v^{-1}(0)$ for $v:M\to N$, $N$
sharp; every face is f.g.\ as a monoid and they are finite in number.
\item A face $K$ is saturated relative to $M$: $K^{\gp}\cap M=K$. The
face generated by $K_1\cup K_2$ has group $K_1^{\gp}+K_2^{\gp}$.
\end{enumerate}
\end{proposition}

\begin{proof}
(a) $x=k_1-k_2\in M$, $k_i\in K$ gives $x+k_2=k_1\in K$, hence
$x\in K$: so $M/K\hookrightarrow M^{\gp}/K^{\gp}$; the rest of the
monoid axioms follow from the face property; initiality is clear.
In $N:=M[-c]^{\sharp}$ let $e$ invert $\bar c$: $\bar c+e=0$ and
sharpness gives $\bar c=0$: hence $N$ receives the universal map from
$M/\langle c\rangle$; conversely killing $c$ inverts it (by $0$): same
universal problem. Poset claim: by initiality.
(b) Kernels into sharp monoids are faces since $v(x)+v(y)=0\Rightarrow
v(x)=v(y)=0$ by sharpness; conversely $K=\ker(M\to M/K)$. A face is a
supporting-hyperplane section of the f.g.\ rational polyhedral cone
$\Cone(M)$: f.g.\ as a monoid (Gordan's lemma), and finitely many in
number.
(c) The first claim is the computation made at the start of (a).
For the second: $\langle K_1\cup K_2\rangle_{\mathrm{face}}$ sits
inside the face $(K_1^{\gp}+K_2^{\gp})\cap M$, whose group is
$K_1^{\gp}+K_2^{\gp}$; conversely any $x\in(K_1^{\gp}+K_2^{\gp})\cap
M$ satisfies $x+k_1+k_2=k_1'+k_2'$ for some $k_i,k_i'\in K_i$ (write
$x$ as a difference and use $x\in M$), so $x$ lies in the face
generated by $K_1\cup K_2$.
\end{proof}

\begin{construction}[Pushouts in $A_P$]\label{constr:pushout}
For $T=(P\to Q_T)$, $X=(P\to Q_X)$, $Y=(P\to Q_Y)$ over $T$: with
$\bar P:=P/\Phi_T$ and $j_X:\bar P\to Q_X$, $j_Y:\bar P\to Q_Y$:
\[
G:=Q_X^{\gp}\oplus_{\bar P^{\gp}}Q_Y^{\gp}
=\bigl(Q_X^{\gp}\times Q_Y^{\gp}\bigr)\big/
\bigl\{(j_X^{\gp}(c),-j_Y^{\gp}(c)) : c\in\bar P^{\gp}\bigr\},
\]
$W_0:=\operatorname{im}(Q_X\times Q_Y\to\breve G)$,
$\breve G:=G/G_{\mathrm{tors}}$, and
$X\otimes_TY:=\bigl(P\to (W_0^{\sat})^{\sharp}\bigr)$. Then
$x\in\Phi_{X\otimes_TY}$ iff $\alpha_X(x)\mapsto0$ in
$(W_0^{\sat})^{\sharp}$.
\end{construction}

\subsection{Finite presentation}\label{ssec:fp}

Birhizomicity at a valuative target includes, by the problem's
definition, the requirement that the map be \emph{of finite
presentation} (f.p.) in the category of monoids. We record what
this means and that every map used as or inside a cover below
satisfies it.

\begin{proposition}\label{prop:fp}
\textup{(}i\textup{)} A monoid map $f:M\to N$ is f.p.\ iff it
presents $N$ from $M$ by finitely many generators and finitely
many relations; equivalently $N$ is a quotient of
$M\oplus\mathbb N^r$ by a finitely generated congruence.
\textup{(}ii\textup{)} For $P$ f.g., a surjective $\varphi:P\to
\mathbb N$ is f.p.: the congruence
$\{(x,y):\varphi(x)=\varphi(y)\}$ is a congruence on the f.g.\
monoid $P\times P$, generated as such by the kernel face
relations and the analogous pointed-cone data; concretely it is
the saturation of a f.g.\ congruence because the graph of
$\ker\varphi^{\gp}\subseteq P^{\gp}\times P^{\gp}$ is a f.g.\
 subgroup and $P$ is saturated \textup{(}this is the standard
f.p.-of-integral-saturated-maps lemma; we give the argument in
the next item\textup{)}.
\textup{(}iii\textup{)} The CFA maps $\pi_K:P\to P/K$ are f.p.\
for every face $K$ of the f.g.\ $P$: $K$ is generated as a face
by finitely many elements $k_1,\dots,k_s$
\textup{(}Proposition~\ref{prop:bp}\textup{(}b\textup{):} f.g.\
as a monoid, and one may take a single element
$c=k_1+\cdots+k_s$\textup{)}, and $P[-c]$ is the localization,
f.p.; the sharpening quotient is generated by the finitely many
unit relations of a f.g.\ monoid, hence f.p.\ again.
\textup{(}iv\textup{)} Pushouts of f.p.\ maps over a common map
of monoids in our category are f.p.: they are quotients of
finite colimits of f.g.\ monoids; saturation and sharpening are
f.p.\ quotients of f.g.\ monoids.
\end{proposition}

\begin{proof}
(i) Standard (Eilenberg-type exercise with congruences).
(iii) $\pi_K$ factors as $P\to P[-K]=P[-c]\to P/K$; each step is
f.p.\ by the displayed presentations: $P[-c]=\langle P,\ e
\mid \bar c+e=0\rangle$-style with one generator one relation;
the sharpening of a f.g.\ monoid $M$ is $M$ modulo the
f.g.-generated superdiagonal congruence $\{(u,-u):u\in
M^*\cup\text{f.g.\ generating set of units}\}$ (the unit group of
a f.g.\ saturated sharp-quotient candidate is f.g.\ because it
is the group of invertible elements of $\Cone(M)$, a finitely
generated rational polyhedral cone).
(ii) $\varphi$ is the composite $P\to P/\ker\varphi
\hookrightarrow\mathbb N$; the first is f.p.\ by (iii); the
second exhibits $\mathbb N$ as $(P/\ker\varphi)$ with the
finitely many relations generated by the image of any finite
generating system of the f.g.\ group $P^{\gp}/\ker\varphi^{\gp}$
--- the congruence determining the halfline inside the rank-one
group is generated by the images of a finite basis of
$\ker\varphi^{\gp}$ (basis finite by f.g.\ of $P^{\gp}$ and
submodule finiteness over $\mathbb Z$): explicitly, choose
$p_0\in P$ with $\varphi(p_0)>0$ minimal in the image group;
then $\mathbb N=(P/\ker\varphi)/\!\sim$ with $\sim$ generated by
$p\sim\varphi(p)\cdot[\bar p_0]$ for $p$ in a generating set of
$P$: f.p.
(iv) The displayed Construction~\ref{constr:pushout} exhibits
$X\otimes_TY$ from f.g.\ monoids by f.g.\ quotients (the
$\bar P^{\gp}$-mixed congruence is generated by the images of a
finite generating set of $\bar P^{\gp}$, f.g.\ when $\bar P$ is
f.g.) and the saturation/sharpening steps of (iii).
\end{proof}

\begin{remark}[Standing hypothesis $P$ f.g.; uses in covers]
\label{rem:nonfg}
From here on, $P$ is finitely generated. Consequently: every
fold $\mathfrak f_\varphi$ and every CFA piece map used in any
cover below is f.p.\ by Proposition~\ref{prop:fp}(ii),(iii); the
pushouts of Construction~\ref{constr:pushout} among them are
f.p.\ by (iv); the test maps $(v/d)$ and
$(d\cdot\mathrm{id}_{\mathbb N})$ of Lemma~\ref{lem:atomic} are
f.p.\ (targets $\mathbb N$); and every refinement map inserted
below (Corollaries~\ref{cor:probecover}, \ref{cor:probecovers},
Lemma~\ref{lem:sandwich}) is f.p. Without f.g.\ of $P$, the CFA
family $\{P\to P/K\}$ may fail to be a cover (faces may fail to
be finitely generated as faces, so $\pi_K$ may fail to be
f.p.), and the face poset can be infinite; these are the two
obstructions to the immediate generalization of the argument to
arbitrary $P$ \textup{(}cf.\ the Scope hypothesis in
\S\ref{sec:intro}\textup{)}.
\end{remark}

\begin{lemma}[Covers may be taken surjective]\label{lem:surj}
Every cover in the slice admits a refinement whose piece maps are
surjective monoid homomorphisms \textup{(}in particular
$Z_U=P+\Phi_U^{\gp}$ for every nerve piece, so
Lemmas~\ref{lem:nerve}, \ref{lem:mixed} apply with the stated
$Z$-hypotheses\textup{)}.
\end{lemma}

\begin{proof}
A fold $\mathfrak f_\varphi$ with non-primitive
$\varphi=d\varphi_0$ \textup{(}$d={GF}(\varphi)$, $\varphi_0$
primitive\textup{)} is refined by the singleton piece
$P\xrightarrow{\varphi_0}\mathbb N$: this covers
$\mathfrak f_\varphi$: a valuative point factoring through
$\mathfrak f_\varphi$ has the form $v=c\varphi_0$ with
$d\mid c$ (both functionals have a common primitive part), and
then $v=(c/d)\circ\varphi_0$; the structure map $\varphi_0$ is
birhizomic at the factor $(c/d\cdot):\mathbb N\to\mathbb N$ by
Lemma~\ref{lem:atomic}, the value groups $(c/d)\Z$ and
$\bigl((c/d)\circ\varphi_0\bigr){}^{\gp}(P^{\gp})=(c/d)\Z$
coinciding since $\varphi_0^{\gp}$ is surjective. The new fold
is surjective. CFA quotient
maps are surjective by construction. Hence every cover admits a
refinement---first the normal form
Theorem~\ref{thm:normalform}, then the primitive reduction of
its folds---whose piece maps are all surjective. For such $U$:
$Z_U=\alpha_U^{\gp,-1}(Q_U)=P+\Phi_U^{\gp}$ \textup{(}any
$q'=\alpha_U^{\gp}(z)$ lifts to $p'\in P$ by surjectivity, and
$z=p'+(z-p')\in P+\ker\alpha_U^{\gp}$ with
$\ker\alpha_U^{\gp}=\Phi_U^{\gp}$:
Lemma~\ref{lem:Zpointed}\textup{)}.
\end{proof}

\subsection{Valuative monoids}

\begin{lemma}\label{lem:N}
Every valuative monoid is $\{0\}$ or $\mathbb N$.
\end{lemma}

\begin{proof}
Let $V\ne\{0\}$ be valuative; $\Gamma:=V^{\gp}$ with convex chain
$0=H_0\subsetneq\cdots\subsetneq H_s=\Gamma$.

\emph{Step 1.} An archimedean ordered group order-embeds into
$(\R,+)$ (H\"older: fix $a>0$; $g\mapsto\sup\{m/n\in\Q:mg\le na\}$);
a f.g.\ subgroup of $\R$ is cyclic or dense.

\emph{Step 2: $s=1$.} If $s\ge2$, fix $\gamma\in H_1$ positive
and $\delta\in H_2\setminus H_1$ positive. For every $n\ge0$,
$\delta-n\gamma\in V$: indeed its image in $\Gamma/H_1$ is
$\bar\delta>0$ in the induced total order. Set
$\Lambda:=H_2/H_1\subseteq\Gamma/H_1$. The ordered group $\Lambda$
has no proper nonzero convex subgroup (the preimage of such a
subgroup in $H_2$ would be convex in $\Gamma$, hence equal to $H_1$
or $H_2$ by maximality of the chain), hence is archimedean; so by
Step 1 it order-embeds in $(\R,+)$. Let $G$ be a finite generating
set of the monoid $V$, and put
$c_0:=\min\{\bar g>0:g\in G,\ \bar g\in\Lambda\}$, where the minimum
over the (finite) set of positive classes exists and is $>0$
because some generator has positive class (otherwise
$V\subseteq H_1$).
Decompose $\delta-n\gamma=\sum_k g_k$ with $g_k\in G$: the
$\Lambda$-classes of the summands are nonnegative with sum
$\bar\delta$, so at most $N_0:=\mu_\Lambda(\bar\delta)/c_0$ of them are
outside $H_1$ (measured through the embedding
$\mu_\Lambda:\Lambda\hookrightarrow\R$). Hence
\[
\delta-n\gamma \;=\; \mathfrak s_n+h_n,\qquad
\mathfrak s_n\text{ a sum of }\le N_0\text{ generators outside
}H_1,\quad h_n\in V\cap H_1 .
\]
The $\mathfrak s_n$ range in a finite set; fix $\mathfrak s$
occurring for infinitely many $n$, and set $w:=\delta-\mathfrak
s\in\Gamma$. Then $w-n\gamma=h_n\in H_1$ for infinitely many $n$;
taking a difference of two such indices shows $w\in H_1$, and
$h_n=w-n\gamma\in V\cap H_1$ means $0\le n\gamma\le w$ in the
archimedean ordered group $H_1$ for infinitely many $n$:
contradiction.


\emph{Step 3: $\Gamma$ not dense.} Suppose $\Gamma$ is dense in
$\R$; put $E:=\Q\Gamma\subseteq\R$, write $\mu_E$ for the inclusion
$E\hookrightarrow\R$, set $E^+:=\{e\in E:\mu_E(e)\ge0\}$ and
$d:=\dim_{\Q}E$, and assume $d\ge2$. Then $V=\Gamma\cap E^+$ and
$C:=\Q_{\ge0}\cdot V$ is a full-dimensional rational polyhedral
cone in $E$ with $C\subseteq E^+$ (full-dimensional because
$\Gamma=V-V$ spans $E$).

\emph{Lemma A: any $\Q$-rational form $\ell$ on $E$ with
$\ell\ge0$ on $V$ is zero.} First, $\mu_E(V)$ is dense in
$\R_{\ge0}$: given $r\ge0$, pick $\delta\in\Gamma$ with
$0<\mu_E(\delta)<1/n$ (dense), and $m\in\N$ with
$|m\mu_E(\delta)-r|<1/n$ (archimedean division in $\R$); then
$m\delta\in V$. Hence, viewing $E_{\R}=\R$ via $\mu_E$ and writing
$\ell_{\R}:=\ell\otimes\R$ (a continuous real form), the
inequality $\ell_{\R}\ge0$ holds on the dense subset
$\mu_E(V)$ of $\R_{\ge0}$, hence on all of $\R_{\ge0}$; therefore
$\ell_{\R}=c\,\mathrm{id}_{\R}$ with $c\ge0$, i.e.\
$\ell(e)=c\,\mu_E(e)$ for all $e\in E$. If $c>0$, the image
$\mu_E(E)=c^{-1}\ell(E)\subseteq c^{-1}\Q$ would be a
$d$-dimensional $\Q$-subspace of $\R$ inside a one-dimensional
one: impossible as $d\ge2$. Hence $c=0$ and $\ell=0$.

\emph{(3b) $C\ne E^+$.} Write $C=\bigcap_{j=1}^{r}
\{\ell_j\ge0\}$ with $\Q$-rational $\ell_j$ (a rational
polyhedral cone in $E$ has rational supporting forms). If
$C=E^+$, each $\ell_j\ge0$ on $E^+\supseteq V$: Lemma A gives
$\ell_j=0$ for all $j$, i.e.\ $C=E$: contradicting
$C\subseteq E^+\subsetneq E$.

\emph{(3c) Conclusion of Step 3.} By (3b) there is
$v\in E^+\setminus C$; strict separation of the closed cone $C$
from the point $v$ inside the $\Q$-vector space (finite
dimensional: standard strict separation applied rationally)
yields a $\Q$-rational form $\ell$ with $\ell|_C\ge0$ and
$\ell(v)<0$. Since $V\subseteq C$, $\ell|_V\ge0$: Lemma A gives
$\ell=0$, contradicting $\ell(v)<0$. Therefore $d=1$: $\Gamma$
is a f.g.\ subgroup of $\Q$, hence cyclic: $\Gamma=\Z\gamma$,
and $V=\N\gamma\cong\N$.
\end{proof}

\begin{corollary}\label{cor:function}
Valuative points are $0$ or functionals $\varphi:M\to\mathbb N$, with
kernel face $\{\varphi=0\}$ and half-space cone
$\{\varphi^{\gp}\ge0\}$. On a sharp $M\ne\{0\}$ tight functionals
exist \textup{(}strictly positive $\Q$-formals on the pointed cone
$\Cone(M)$, denominator-cleared\textup{)}.
\end{corollary}

\subsection{Infinitesimal maps}

\begin{lemma}\label{lem:inj0}
If $p:M\to N$ is infinitesimal with $N$ sharp then $\ker p=\{0\}$; if
$N$ is moreover valuative then $p^{\gp}$ is injective.
\end{lemma}

\begin{proof}
As usual $x\prec y$ means $my\le x\le ny$ for some $m,n\in\Z$.
First: let $x\in M$, $p(x)=0$; with $X:=-x$, $y:=x$: $X\prec y$;
$p(X)=0\in N$; $p(y)\prec0$ reads $0\le p(y)\le0$ \checkmark; the condition
gives $-x\in M$, a unit: $x=0$.
Second: $X=x_1-x_2\in\ker p^{\gp}$; $y:=x_1+x_2$: $X\prec y$;
$p(X)=0$; $p(y)\prec0$ gives $p(x_1)+p(x_2)=0$ in $N$; $N$ sharp:
each $=0$; first part: $x_i=0$.
\end{proof}

\begin{lemma}\label{lem:infstruct}
For $p:M\to\mathbb N$ nonzero: $p$ infinitesimal iff $p^{\gp}$
injective and $M$ conical: $M=\{x\in M^{\gp}:p^{\gp}(x)\ge0\}$; then
$\pi:=p/{GF}(p):M\cong\mathbb N$.
\end{lemma}

\begin{proof}
``If'' is immediate from the premise $p(x)\in\mathbb N$ alone. ``Only
if'': $\ker p^{\gp}=0$ by Lemma~\ref{lem:inj0}; for conicality let
$p^{\gp}(x)>0$, write $x=x_1-x_2$, $y:=x_1+x_2$: $x\prec y$;
$p(y)\prec p(x)$: $\mathbb N$ is archimedean: $p(y)\le kp(x)$ for
some $k$: premises OK: $x\in M$.
\end{proof}

\subsection{Atomicity}

\begin{lemma}[Atomicity]\label{lem:atomic}
Let $f:Q\to Q'$ be f.p.\ \textup{(}this is part of birhizomicity;
all our applications have it by
Proposition~\ref{prop:fp}\textup{)} and
$\psi:Q'\to\mathbb N$ nonzero. Then $f$
is birhizomic at $\psi$ iff
$\psi^{\gp}(Q'^{\gp})=(\psi f)^{\gp}(Q^{\gp})$.
\end{lemma}

\begin{proof}
\emph{Necessity.} With $v:=\psi f$, $d:={GF}(v)$, test the pair
$(h,p):=(v/d,\ d\cdot\mathrm{id}_{\mathbb N})$
\textup{(}infinitesimal by Lemma~\ref{lem:infstruct}; both maps
are f.p.\ by Proposition~\ref{prop:fp}(ii)/(iv), and the source
$M=\mathbb N$ is f.g.\textup{)}; a lift $u$ satisfies $du=\psi$, hence
$\psi^{\gp}(Q'^{\gp})\subseteq d\Z=(\psi f)^{\gp}(Q^{\gp})$; the
other inclusion is automatic.
\emph{Sufficiency.} Given $(h,p)$: ${GF}(p)$ divides ${GF}(v)$
($v^{\gp}=p^{\gp}h^{\gp}$), so $\psi^{\gp}(Q'^{\gp})
={GF}(v)\Z\subseteq{GF}(p)\Z$; define
$u^{\gp}:=(\pi^{\gp})^{-1}\circ(\psi^{\gp}/{GF}(p))$ where
$\pi:=p/{GF}(p)$ (isomorphism, Lemma~\ref{lem:infstruct});
$u^{\gp}$ maps $Q'$ into $M$ by conicality and
$p^{\gp}u^{\gp}=\psi^{\gp}$ (so $pu=\psi$ on $Q'$: on the monoid:
$p^{\gp}u^{\gp}(q')=\psi^{\gp}(q')\in\mathbb N$ \checkmark); and
$uf=h$: $p^{\gp}u^{\gp}f^{\gp}=\psi^{\gp}f^{\gp}=p^{\gp}h^{\gp}$,
$p^{\gp}$ injective; uniqueness: same injectivity.
\end{proof}

\begin{corollary}[CFA pieces and covers]\label{cor:cfa}
Let $Q$ be a monoid and $K$ a face.
\textup{(}a\textup{)} $Q\to Q/K$ is birhizomic at every
$g:Q/K\to V$, $V$ valuative.
\textup{(}b\textup{)} $\{Q\to Q/K:K\text{ face of }Q\}$ is a
birhizomic cover \textup{(}the \emph{CFA cover}\textup{)}.
\end{corollary}

\begin{proof}
(a) Nonzero $g=\psi$: Lemma~\ref{lem:atomic}: value-group equality
holds since $\pi_K^{\gp}$ is surjective. Zero $g$: $ph=0$ and
Lemma~\ref{lem:inj0} give $h=0$; unique lift.
(b) At $v$: $v=0$: piece $K=Q$ with (a); $v=\varphi\ne0$: piece
$K=\ker\varphi$ with $\bar\varphi$ and (a).
\end{proof}

\begin{theorem}[Normal form]\label{thm:normalform}
Let $T=(P\to Q_T)$. Every birhizomic cover of $T$ admits a refinement
whose pieces are fold pieces $\mathfrak f_\varphi=
(P\xrightarrow{\varphi}\mathbb N)$ with $\varphi(\Phi_T)=0$, and CFA
pieces $(P\to P/K)$ with $K\supseteq\Phi_T$.
\end{theorem}

\begin{proof}
For each piece $S=(P\xrightarrow{\alpha}Q)$ of the cover, the family
\[
\mathcal N(S)=
\{\mathfrak f_{\theta\alpha}:\theta:Q\to\mathbb N\}
\ \cup\ \{Q/\mathfrak k:\mathfrak k\text{ face of }Q\}
\]
is a cover of $S$: indeed, let $v:P\to V$ factor through
$\alpha$ (these are the points covers of $S$ are responsible for).
If $V=\{0\}$: use $\mathfrak k=Q$. If $v=\theta\alpha$,
$\theta:Q\to\mathbb N$ nonzero: write $\theta=d\theta_1$ with
${GF}(\theta_1)=1$; then $v=(d\cdot)\circ(\theta_1\alpha)$ through
the fold piece $\mathfrak f_{\theta_1\alpha}$; birhizomicity of the
structure map $\theta_1$ at $g=d\cdot\mathrm{id}$:
Lemma~\ref{lem:atomic}: the condition
$g^{\gp}(\Z)=(g\theta_1)^{\gp}(Q^{\gp})$ reads
$d\Z=\theta^{\gp}(Q^{\gp})=d\Z$ \checkmark. Compose these covers over all
pieces (transitivity).
\end{proof}

\begin{remark}\label{rem:cover-points}
In the slice topology, covers of $T$ only need to answer those
$v:P\to V$ that factor through $\alpha_T:P\to Q_T$ --- in
particular $v(\Phi_T)=0$ and $v$ factors through
$P\to P/\Phi_T$: the problem's definition applied to the object
$\alpha_T$ quantifies over $v=Q_T\to V$, and such $v$ kill the
kernel of $\alpha_T$; all our covering constructions
(Theorem~\ref{thm:normalform}, Corollaries~\ref{cor:cfa},
\ref{cor:probecover}, \ref{cor:probecovers}) answer such $v$ by
structure maps killing $\Phi_T$ \emph{because} the pieces used
(folds with $\varphi(\Phi_T)=0$, quotients $P/K$ with
$K\supseteq\Phi_T$) receive $\alpha_T$ (they lie over $T$).
\end{remark}


\subsection{Pushout kernels against folds}

Fix $T$, $J:=\Phi_T$, $\bar P:=P/J$.

\begin{lemma}\label{lem:Zpointed}
For $U$ over $T$: $\alpha_U^{\gp}(L_U)=\{0\}$ and
$P\cap L_U=\Phi_U$; the induced
$Z_U/L_U\hookrightarrow Q_U$ is injective with image the sharp
saturated submonoid $Q_U\cap\alpha_U^{\gp}(P^{\gp})$; the group
$J^{\gp}$ satisfies $J^{\gp}\subseteq\ker\alpha_U^{\gp}$.
\end{lemma}

\begin{proof}
$\lambda\in L_U$: $\pm\alpha_U^{\gp}(\lambda)\in Q_U$; sharp:
$=0$. Hence $x\in P\cap L_U\Rightarrow\alpha_U(x)=0$:
$P\cap L_U\subseteq\Phi_U$; conversely $\Phi_U\subseteq L_U$ directly.
Injectivity of $Z_U/L_U\to Q_U$: $\alpha_U(z_1-z_2)=0$ with
$z_i\in Z_U$: then $\pm(z_1-z_2)\in Z_U$: $z_1-z_2\in L_U$. For $J$:
$\alpha_U(J)=0$.
\end{proof}

\begin{lemma}[Kernel criterion]\label{lem:pushout}
Let $U$ over $T$ and $\mathfrak f:=\mathfrak f_\varphi$ with
$\varphi(J)=0$; $W:=U\otimes_T\mathfrak f$. For $x\in P$:
\[
x\in\Phi_W\ \Longleftrightarrow\ x\in\Phi_U\ \ \text{or}\ \
\exists\,t\ge1,\,a\in P^{\gp}:\ a-tx\in Z_U,\ \varphi^{\gp}(a)\le0 .
\]
\end{lemma}

\begin{proof}
Reduce modulo $J$; write $j:\bar P\to Q_U$. By
Construction~\ref{constr:pushout}, $\bar x$ dies in $W$ iff there
exist $t\ge1$ and a unit $u$ of $W_0^{\sat}$ with
$[(\alpha_U(tx),0)]+u=0$ in $G$.

\emph{Necessity: the pair (A)+(B).} Suppose $x\in\Phi_W$; if
$x\in\Phi_U$ there is nothing to add. Otherwise take $(t,u)$ as
above. Writing the equality out with a class $u=[(\xi,m)]$: there
is $c\in\bar P^{\gp}$ with $\alpha_U(tx)+\xi=j(c)$ and
$m=-\varphi^{\gp}(c)$; lift $c$ to $a_0\in P^{\gp}$ and set
$z_0:=a_0-tx$: then $\alpha_U^{\gp}(z_0)=\xi$. The condition
$+u\in W_0^{\sat}$ means that a positive multiple $s u$ has a
representative in $Q_U\times\mathbb N$; raw representatives of
$s u$ are $(\alpha_U^{\gp}(s z_0),\,-s\varphi^{\gp}(a_0))$,
measured up to slides $(\xi,m)\mapsto(\xi+j(d),
m-\varphi^{\gp}(d))$, $d\in\bar P^{\gp}$: hence for some $d$:
$\alpha_U^{\gp}(s z_0+d)\in Q_U$ and
$s\varphi^{\gp}(a_0)+\varphi^{\gp}(d)\le0$. Take
$(t',a):=(s t,\ s a_0+\hat d)$ with $\hat d$ lifting $d$; then
$a-t'x=s z_0+\hat d\in Z_U$ \textup{(}condition
\textup{(A)}\textup{)} and $\varphi^{\gp}(a)\le0$
\textup{(}condition \textup{(B)}\textup{)}: the required pair
exists.

\emph{Sufficiency.} Conversely, given $(t,a)$ with (A),(B): set
$z:=a-tx$, $\xi:=\alpha_U^{\gp}(z)\in Q_U$,
$m:=-\varphi^{\gp}(a)\ge0$, and $u:=[(\xi,m)]\in W_0$; the
relation $a=tx+z$ gives
$[(\alpha_U(tx),0)]+u=[(j(\bar a),-\varphi^{\gp}(\bar a))]=0$ in
$G$; hence $[(\alpha_U(tx),0)]=-u$ lies in $W_0$ as well, so $u$
is a unit, $t\bar x$ dies after sharpening, and $x$ dies
(saturation): $x\in\Phi_W$. That
$x\in\Phi_U\Rightarrow x\in\Phi_W$ is clear (the structure map of
$W$ factors through that of $U$). \end{proof}

\begin{example}[Fold vs fold]\label{ex:ff}
$U=\mathfrak f_\psi$, $\psi$ tight on $\bar P$, and
$\mathfrak f_\varphi$: the criterion (with
$Z_U=\{\psi^{\gp}\ge0\}$): if $[\varphi]=[\psi]$:
(B)$\Rightarrow\psi(a)\le0$, (A): $\psi(a)\ge t\psi(x)\ge1$ for
$x\ne0$: impossible: $\Phi_W=J$. If $[\varphi]\ne[\psi]$: the open
rational wedge $\{\varphi\le0<\psi\}$ in $\bar P^{\gp}\otimes\R$ is
nonempty; scale $a$ in it with $\psi(a)\ge t\psi(x)$: every
$x\neq0$ in $\bar P$ dies: $\Phi_W=\bar P$.
\end{example}

\begin{proposition}[Equiface probes]\label{prop:equiface}
Let $U$ over $T$ and $\varphi:P\to\mathbb N$ satisfy
\textup{(}G\textup{)}: $\varphi(L_U)=0$ and $\varphi(z)>0$ for
$z\in Z_U\setminus L_U$ \textup{(}such $\varphi$ exist, and may be
taken \emph{primitive}: $Z_U/L_U\hookrightarrow Q_U$ is a sharp f.g.\
monoid by Lemma~\textup{\ref{lem:Zpointed}}, so a strictly positive
functional on it exists by Corollary~\textup{\ref{cor:function}};
dividing out the content preserves the vanishing and positivity
pattern; then automatically
$\varphi(J)=0$ since $J^{\gp}\subseteq L_U$\textup{)}. Then:
\textup{(}i\textup{)} $\Phi_{\mathfrak f_\varphi}=\Phi_U$;
\textup{(}ii\textup{)} $\Phi_{U\otimes_T\mathfrak f_\varphi}=\Phi_U$.
In particular, for the base piece $iB(T)$ \textup{(}where
$Z=P+J^{\gp}$, $L=J^{\gp}$\textup{)}, \emph{every} tight functional
$\varphi$ of $\bar P$ satisfies \textup{(}G\textup{)}, so for every
tight fold $\mathfrak f$:
$\Phi_{\mathfrak f}=J$ and
$\Phi_{\mathfrak f\otimes_T iB(T)}=J$.
\end{proposition}

\begin{proof}
(i) $\Phi_{\mathfrak f_\varphi}=\{x\in P:\varphi(x)=0\}$: $x\in P
\subseteq Z_U$ with $\varphi$-zero: (G): $x\in L_U$: $x\in
P\cap L_U=\Phi_U$ (Lemma~\ref{lem:Zpointed}) \checkmark.
(ii) Lemma~\ref{lem:pushout}: suppose $x\notin\Phi_U$ dies: $(t,a)$:
$z:=a-tx\in Z_U$, $\varphi(a)\le0$. If $\varphi(z)<0$: but $z\in
Z_U$ and (G) gives $\varphi(z)\ge0$ (with $>0$ off $L_U$):
contradiction. So $\varphi(z)\ge0$; then
$\varphi(z)=\varphi(a)-t\varphi(x)$ gives
$\varphi(a)\ge t\varphi(x)$; with $\varphi(a)\le 0$ and
$\varphi(x)>0$ (as $x\in P\subseteq Z_U$, $x\notin L_U$ since
$x\notin\Phi_U=P\cap L_U$): $0\ge\varphi(a)\ge t\varphi(x)\ge1$:
contradiction. Hence $\Phi_W=\Phi_U$ \checkmark.
The $iB(T)$-computation: $Z_{iB(T)}=P+J^{\gp}$,
$L_{iB(T)}=J^{\gp}$ \checkmark (sharpness of $\bar P$: $\pm\bar x\in\bar P\Rightarrow
\bar x=0$, i.e.\ $x\in J$); (G) for tight $\varphi$:
$\varphi(J^{\gp})=0$ \checkmark; $z\in Z\setminus L$ has $\bar z\ne0$, so
$\varphi(\bar z)>0$ \checkmark.
\end{proof}

\begin{corollary}[Probes exist in covers]\label{cor:probecover}
Let $R=\{T\to T_i\}$ be a cover of $T$, and let
$\varepsilon:P\to\mathbb N$ be any functional \emph{coming from an
equiface probe}, i.e.\ satisfying for some piece $U$ over $T$ the
condition \textup{(}G\textup{)} of
Proposition~\textup{\ref{prop:equiface}}
\textup{(}in particular $\Phi_{\mathfrak f_\varepsilon}=\Phi_T$
by \textup{(}i\textup{)} loc.\ cit.\textup{)}. Then $\varepsilon$ may be
taken primitive \textup{(}replace by its primitive
part\textup{)}, and $R$ admits a refinement containing
$\mathfrak f_\varepsilon$. In particular
\textup{(}with $U=iB(T)$, where every tight $\varepsilon$
works\textup{)}: every cover admits a refinement containing a tight
fold with kernel exactly $\Phi_T$.
\end{corollary}

\begin{proof}
By Proposition~\ref{prop:equiface}, equiface probes may be
chosen primitive: replace the given $\varepsilon$ by its
primitive part, which satisfies the same vanishing and
positivity conditions, hence also defines a fold with kernel
$\Phi_T$. Apply the cover property of $R$ at the point
$v:=\varepsilon$: some member $T_i$ admits
$\theta:Q_{T_i}\to\mathbb N$ with $\theta\alpha_i=v$ and
$T\to T_i$ birhizomic at $\theta$; by
Lemma~\ref{lem:atomic},
$\theta^{\gp}(Q_{T_i}^{\gp})=v^{\gp}(P^{\gp})$, i.e.\
${GF}(\theta)={GF}(v)=1$. Refine $T_i$ by its CFA cover
(Corollary~\ref{cor:cfa}(b)) and keep the member over the face
$\ker\theta$: the induced $\bar\theta:Q_{T_i}/\ker\theta
\hookrightarrow\mathbb N$ is injective with saturated sharp
${GF}{=}1$-image in $\mathbb N$, hence an isomorphism; the
member $S:=(P\to Q_{T_i}/\ker\theta)$ is therefore isomorphic
over $P$ to $\mathfrak f_\varepsilon$.
\textup{(}All maps used are f.p.: the slice maps by
Proposition~\ref{prop:fp}(ii)--(iii), and the refinement
composites by (iv) therein.\textup{)}
The composed refinement of
$R$ contains $\mathfrak f_\varepsilon$.
\end{proof}

\begin{corollary}[Probe covers]\label{cor:probecovers}
Let $\bar P\ne\{0\}$ and let $\chi:\bar P\to\mathbb N$ be a primitive
tight functional. Then the family
\[
\square_1\;:=\;\{\mathfrak f_{\chi\bar\pi}\}\ \cup\
\{(P\to P/K):K\supseteq J\text{ face}\}
\]
\textup{(}one tight primitive fold on the ray $[\chi]$, together with
all CFA pieces\textup{)} already covers $T$.
\end{corollary}

\begin{proof}
Let $v=\psi\bar\pi$, $\psi:\bar P\to\mathbb N$ nonzero;
$\ker\psi=:K^0\supseteq0$ a face of $\bar P$; let
$K:=$ preimage of $K^0$ in $P$ (a face $\supseteq J$). Use the piece
$P/K$: $v$ factors $P/K\xrightarrow{\bar\psi}\mathbb N$;
birhizomicity of the structure map $P\to P/K$ at $\bar\psi$:
Corollary~\ref{cor:cfa}(a). If $\ker\psi=0$ and $[\psi]=[\chi]$:
$\psi=d\chi$: use the fold piece and Lemma~\ref{lem:atomic}:
requires $d\Z=(\psi\bar\pi)^{\gp}(P^{\gp})=d\Z$ \checkmark. The zero point:
dead piece.
\end{proof}


\section{The inverse image is a sheaf}\label{sec:sheaf}

\begin{proposition}\label{prop:adj}
$i:B_P\to A_P$ has right adjoint $\varphi_0:A_P\to B_P$,
$T\mapsto B(T)$; hence the pointwise left Kan extension of $F$ along
$i$ is $G_F:=F\circ\varphi_0$, and $G_F(i(b))=F(b)$.
\end{proposition}

\begin{proof}
$\operatorname{Hom}_{A_P}(i(P\to P/K),T)$ is a singleton iff
$K\subseteq\Phi_T$, else empty (Proposition~\ref{prop:bp}), matching
$\operatorname{Hom}_{B_P}((P\to P/K),B(T))$: adjunction. Kan
extension along a functor with a right adjoint is precomposition with
the adjoint; pointwise: $(i\downarrow T)$ has terminal object $B(T)$.
\end{proof}

\begin{proposition}[Separatedness]\label{prop:sep}
For every cover $R$ of $T$, $G_F(T)\to\prod_{T_i\in
R}G_F(T_i)$ is injective.
\end{proposition}

\begin{proof}
By Theorem~\ref{thm:normalform} and Corollary~\ref{cor:probecover},
$R$ admits a refinement (pulling back along which the map in question
factors) containing a tight fold $\mathfrak f$ with
$\Phi_{\mathfrak f}=J$; the restriction $G_F(T)\to
G_F(\mathfrak f)$ is $F(\mathrm{id}_{B(T)})$: injective.
\end{proof}

\begin{theorem}[Descent]\label{thm:descent}
For every cover $R$ of $T$, $G_F(T)\xrightarrow{\sim}\HH^0(R;G_F)$.
Hence $G_F$ is a sheaf on $(A_P^{\op})_{\mathrm{birhiz}}$;
consequently $\pi^{-1}F=G_F$ and
$F(b)\xrightarrow{\sim}(\pi_*\pi^{-1}F)(b)$, canonically.
\end{theorem}

\begin{proof}
Injectivity is Proposition~\ref{prop:sep}. Let
$(s_U)\in\HH^0(R;G_F)$. By Theorem~\ref{thm:normalform},
Corollary~\ref{cor:probecover}, and
Corollary~\ref{cor:probecovers}, we may refine (and pull the system
back; this does not change its fate for surjectivity) so that $R$
contains: a tight fold $\mathfrak f_\star$; the base piece
$iB:=i(B(T))$; all CFA pieces $(P\to P/K)$, $K\supseteq J$; and for
every piece $U$ of $R$ a (G)-probe $\mathfrak f_{\varepsilon_U}$
with $\varepsilon_U$ satisfying (G) for $U$
(Proposition~\ref{prop:equiface}). The nerve of the refined cover
contains all pushouts of its pieces. Set
\[
s^0\;:=\;s_{iB}\in F(B(T)).
\]
\emph{Claim A: $s_\star=s^0$.} The kernel of
$W:=\mathfrak f_\star\otimes_T iB$ is $J$
(Proposition~\ref{prop:equiface}, $iB$-case: every tight functional
satisfies (G) for $iB$). Compatibility at $W$: the two restriction
maps $G_F(\mathfrak f_\star)\to G_F(W)$ and $G_F(iB)\to G_F(W)$ are
both $F(\mathrm{id}_{B(T)})$ (as $\Phi_\star=\Phi_W=\Phi_{iB}=J$);
hence $s_\star=s_{iB}=s^0$.

\emph{Claim B: for every piece $U$ of $R$,
$s_U=F(B(T)\twoheadrightarrow B(U))(s^0)$.}
\emph{Case $\Phi_U=J$.} The probe piece $\mathfrak
f_{\varepsilon_U}$ has kernel $J$ (Prop.~\ref{prop:equiface}(i)),
i.e.\ $\varepsilon_U$ is tight on $\bar P$, hence satisfies (G) for
$iB$ (Prop.~\ref{prop:equiface}, last clause); so
$\Phi_{\mathfrak f_{\varepsilon_U}\otimes_T iB}=J$ and the two
restriction maps into the overlap are
$F(\mathrm{id}_{B(T)})$: compatibility gives
$s_{\varepsilon_U}=s_{iB}=s^0$. Also, by
Prop.~\ref{prop:equiface}(ii),
$U\otimes_T\mathfrak f_{\varepsilon_U}$ has kernel $\Phi_U=J$
with identity transitions from both factors, so
$s_U=s_{\varepsilon_U}=s^0$, the required form since
$B(U)=B(T)$.
\emph{Case $\Phi_U=K\supsetneq J$.} The probe piece
$\mathfrak f_{\varepsilon_U}$ has kernel $K$
(Prop.~\ref{prop:equiface}(i)), so
$U\otimes\mathfrak f_{\varepsilon_U}$ has kernel $K$
(loc.\ cit.\ (ii)) with identity transitions:
$s_U=s_{\varepsilon_U}$. We now compute the kernel of
$W:=\mathfrak f_{\varepsilon_U}\otimes_T iB$ via
Lemma~\ref{lem:pushout} ($Z_{iB}=P+J^{\gp}$; write bars for
reduction mod $J$, so $Z_{iB}$-membership is
``$\overline{\,\cdot\,}\in\bar P$'', and
$\varepsilon_U=\psi\bar\pi$ with $\psi:\bar P\to\mathbb N$ whose
kernel is exactly the face $K/J$):
\emph{claim:} $\Phi_W=K$.

If $\bar x\in K/J$: take $t:=1$ and a lift $a$ of $\bar a:=\bar x$;
then $a-x$ reduces to $0$, hence lies in $J^{\gp}\subseteq
P+J^{\gp}=Z_{iB}$: (A) holds; and
$\varepsilon_U(a)=\psi(\bar x)=0$: (B) holds. Hence
$x\in\Phi_W$, i.e.\ $K\subseteq\Phi_W$.

Conversely let $x\in\Phi_W$ via $(t,a)$: (B): $\psi(\bar a)\le0$,
with $\bar a=\overline{a-tx}+t\bar x$, where
$\bar z:=\overline{a-tx}\in\bar P$ by (A). Two cases: if
$\psi(\bar a)=0$, then $\bar a\in K/J$; the equation
$t\bar x=\bar a-\bar z$ with $\bar z\in\bar P$ is an equality
$t\bar x+\bar z=\bar a$ inside $\bar P$, i.e.\ $t\bar x$ and
$\bar z$ are elements of $\bar P$ whose sum lies in the face
$K/J\subseteq\bar P$; by the face property, $t\bar x\in K/J$,
whence $\bar x\in K/J$ (faces are saturated;
Proposition~\ref{prop:bp}(c)), i.e.\ $x\in K$. If
$\psi(\bar a)<0$: then
$t\psi(\bar x)=\psi(\bar a)-\psi(\bar z)\le\psi(\bar a)<0$ while
$\psi(\bar x)\ge0$: impossible. Hence
$\Phi_W\subseteq K$: equality.
Compatibility at $W$: $G_F(\mathfrak f_{\varepsilon_U})\to G_F(W)$
is $F(\mathrm{id}_{P/K})$ while $G_F(iB)\to G_F(W)$ is
$F(B(T)\twoheadrightarrow P/K)$; hence
$s_{\varepsilon_U}=F(B(T)\twoheadrightarrow P/K)(s^0)$, and with
$s_U=s_{\varepsilon_U}$ above, Claim B follows.
\emph{Endgame.} Claim A+B: the pulled system on the final refinement
$R'$ has the form $\iota_{R'}(s^0)$ with $s^0\in F(B(T))=G_F(T)$: it
lies in the image of $G_F(T)$. Hence every class of
$\colim_R\HH^0(R;G_F)=G_F^+(T)$ lies in the image of $G_F(T)$, so
$G_F(T)\to G_F^+(T)$ (injective by Proposition~\ref{prop:sep}) is an
isomorphism. For the original cover $R$, consider
\[
G_F(T)\longrightarrow\HH^0(R;G_F)\longrightarrow G_F^+(T):
\]
the composite is bijective and the first map is surjective (a
constant system is compatible), so the second map is injective; both
are therefore bijective: $\HH^0(R;G_F)=G_F(T)$ on the nose.
$G_F$ is then a sheaf (descent on all covers: glue maps iso),
$\pi^{-1}F=G_F$, and $(\pi_*\pi^{-1}F)(b)=G_F(i(b))=F(b)$
(Proposition~\ref{prop:adj}).
\end{proof}



\section{Higher direct images}\label{sec:direct}

Let $b=(P\to P/J)$, $\bar P=P/J$, $\mathcal G:=G_F=\pi^{-1}F$ (a
sheaf by Theorem~\ref{thm:descent}). We now compute
$H^q(i(b),\mathcal G)$.

\begin{definition}[The \v{C}ech nerve of a cover]\label{def:nerve}
Let $R=\{i(b)\to T_i\}$ be a cover of $X:=i(b)$ in the sliced site.
Its \emph{nerve} is the simplicial object of the slice whose degree
$q$ part is the disjoint family of the iterated pushouts
$\overline\tau:=\overline{U_0}\cap\cdots\cap
\overline{U_q}:=U_0\otimes_{i(b)}\cdots\otimes_{i(b)}U_q$
over alternating tuples $\tau=(U_0,\dots,U_q)$ of pieces
\textup{(}the pushouts exist and their piece maps are finitely
presented, by Construction~\ref{constr:pushout} and
Proposition~\ref{prop:fp}\textup{(}iv\textup{)}\textup{)}; face
maps are given by forgetting a factor, degeneracies by repeating
one; the alternating \v{C}ech complex of a presheaf $\mathscr W$
is
\[
\cC^{q}(R;\mathscr W)\;=\!\!\prod_{\tau\text{ alternating}}
\mathscr W(\overline\tau),\qquad
(\partial e)_\tau=\sum_{i=0}^{q+1}(-1)^i\,
\mathscr W(d_i)\bigl(e_{\partial_i\tau}\bigr) ,
\]
where $d_i:\overline\tau\to\overline{\partial_i\tau}$ is the
structure map forgetting the $i$-th factor.
\end{definition}

\begin{lemma}[Structure maps between nerve pieces]\label{lem:nervemaps}
For nerve pieces $W_1,W_2$ over $i(b)$ of covers whose pieces
are folds or CFA-quotients \textup{(}as in all covers used
below\textup{)}: a structure map $W_1\to W_2$ over $X$ exists iff
$\Phi_{W_1}\subseteq\Phi_{W_2}$ provided both structure maps are
surjective; in general the surjective pieces among the nerve
\textup{(}$iB$, the $i(P/K)$, the folds satisfied
\textup{)} satisfy the criterion, and the kernel-equal case gives
the identity transition.
\end{lemma}

\begin{proof}
If the structure map of $W_2$ is a face quotient $P/K_2$ or the
fold $\mathbb N$-projection, a map into it out of $W_1$ is the
same as a factorization killing the kernel (universal property
of $P/K_2$: Proposition~\ref{prop:bp}(a); for the folds, maps
between folds on the same ray are the coordinate projections).
Killing $\Phi_{W_2}$ in $W_1$ is possible exactly when
$\Phi_{W_1}\subseteq\Phi_{W_2}$. When the kernels agree, the
induced map on $B$-levels is the identity of
$P/\Phi$, i.e.\ $F(\mathrm{id})$.
\end{proof}

\begin{remark}[The comparison theorem used]\label{rem:verdier}
We use Verdier's hypercovering theorem (SGA~4,
Exp.~V, Th\'eor\`eme~7.4.1), in the following consequence:
on a site with fibre products --- the sliced birhizomic site has
them by Construction~\ref{constr:pushout} --- the sheaf
cohomology of an abelian sheaf $\mathscr W$ on $X$ is computed
by
$H^q(X,\mathscr W)=\colim_{R}\HH^{q}(R;\mathscr W)$,
the colimit over all covers of the alternating \v{C}ech
cohomology of their nerves as in Definition~\ref{def:nerve}; and
a class vanishes in $H^q(X,\mathscr W)$ iff it pulls to zero
along some refinement. (No other external input is used.)
\end{remark}

\subsection{The nerve complex: depth, strands, and transfers}
\label{ssec:nervetricks}

Throughout, $R$ is a cover of $X$ in the sliced site, and all
presheaves take values in abelian groups. For an alternating tuple
$\tau=(U_0,\dots,U_q)$ of pieces write
$\overline\tau:=U_0\otimes_X\cdots\otimes_X U_q$ for its nerve
piece (Definition~\ref{def:nerve}), and
$\partial_i\tau:=(U_0,\dots,\widehat{U_i},\dots,U_q)$.

\begin{remark}[Normalization]\label{rem:normalize}
Definition~\ref{def:nerve} defined $\cC^\bullet(R;\mathscr W)$
directly on \emph{alternating} tuples. Equivalently one may start
from the unreduced complex over all tuples and pass to the quotient
by the degenerate subcomplex: the two compute the same cohomology
(the degenerate cochains form an acyclic subcomplex: standard
normalization of simplicial abelian groups, used here only as this
identification). We work throughout with alternating tuples.
\end{remark}

\begin{definition}[$\mathscr W$-strand of a family]\label{def:strand}
For a subfamily $\mathcal J$ of pieces, the \emph{$\mathcal
J$-strand} is the subquotient of $\cC^\bullet$ obtained by
restricting to tuples all of whose pieces lie in $\mathcal J$: it is
a subcomplex when it is closed under $\partial$ (i.e.\ whenever
every face of a tuple in $\mathcal J$, with its values transported
along the structure map, still has its entries in $\mathcal J$);
more generally, a family of strands is \emph{differential-closed}
when the differential of a cochain supported on it stays supported
on it.
\end{definition}

\begin{lemma}[Constant complexes]\label{lem:constcx}
Let $\mathcal J$ be a differential-closed family of pieces on
which $\mathscr W$ takes a constant value $A$ with all structure
maps inducing $\mathrm{id}_A$, and suppose $\mathcal J$ has a
minimum piece $H$ \textup{(}$H\subseteq U$, hence
$\overline H\hookrightarrow\overline U$ and the induced
$\mathscr W$-map is $\mathrm{id}_A$, for all
$U\in\mathcal J$\textup{)} so that every nerve piece of
$\mathcal J$ receives a structure map from $\overline H$.
Then the $\mathcal J$-strand computes $\HH^0=A$ and
$\HH^{q\ge1}=0$.
\end{lemma}

\begin{proof}
For $q\ge1$ define $(he)_\tau:=e_{(H,\tau)}$ at alternating
tuples $\tau$ of $\mathcal J$ \textup{(}when $(H,\tau)$ has a
repetition its entry is $0$, consistent with the alternating
convention\textup{)}. Since
$\overline{(H,\tau)}=\overline H\otimes\overline\tau$ receives
$\overline H$, the value $e_{(H,\tau)}$ is read off along the
structure map out of the $H$-level, i.e.\ is an element of
$A$. The face relations
$\partial_0(H,\tau)=\tau$,
$\partial_{i+1}(H,\tau)=(H,\partial_i\tau)$ give, with all
structure maps inducing $\mathrm{id}_A$,
\[
(\partial h e)_\tau=\sum_{i=0}^{q}(-1)^i e_{(H,\partial_i\tau)},
\qquad
(h\partial e)_\tau=e_\tau-\sum_{i=0}^{q}(-1)^i
e_{(H,\partial_i\tau)},
\]
so $(\partial h+h\partial)e=e$: $\HH^{q\ge1}=0$. For $q=0$: the
cocycle condition at the pair $(H,U)$ reads
$e_U=\mathscr W(\overline H\hookrightarrow\overline U)(e_H)=e_H$,
so $\HH^0=A$ by evaluation at $H$.
\end{proof}

\begin{lemma}[Exact subquotient and signed transport]\label{lem:subquot}
Let $\cC^\bullet=\cC'\oplus\cC''$ be a decomposition into
differential-closed subcomplexes with $\cC'$ acyclic. Then
$\HH^\bullet(\cC)=\HH^\bullet(\cC'')$; and if
$z=(z',z'')$ is a degree-$q\ge1$ cocycle, there is a cochain
$w'\in\cC'^{\,q-1}$ with $z'=\partial w'$, so that
$z-\partial w'$ is a cocycle supported on $\cC''$ whose
$\cC''$-component differs from $z''$ by the signed transport
$z''-(\mathrm{pr}''\circ\partial)w'$ \textup{(}this component is a
cocycle; if $\cC''$ is closed under $\partial$ on its own it equals
$z''$\textup{)}.
\end{lemma}

\begin{proof}
Acyclicity of $\cC'$ gives $w'$ with $z'=\partial w'$ (the cocycle
$z$ restricts to a cocycle on $\cC'$ by differential-closedness).
Then $z-\partial w'=(0,\,z''-\mathrm{pr}''\partial w')$: a cocycle
supported on $\cC''$, as stated; its cohomology class equals that
of $z$, and the long exact sequence of
$0\to\cC''\to\cC\to\cC'\to0$ identifies $\HH^\bullet(\cC)$
with $\HH^\bullet(\cC'')$ whenever $\cC'$ is acyclic.
\end{proof}

\begin{lemma}[Depth of nerve pieces]\label{lem:depth}
Assign to each piece $U$ over $X$ its \emph{depth}: the length of
a maximal chain in the face poset of $P$ from $J$ up to the kernel
face $\Phi_U$ \textup{(}well-defined since the face poset is
finite, Proposition~\ref{prop:bp}\textup{(}b\textup{)}; the actual
constants never matter\textup{)}; the dead piece $D$
\textup{(}kernel $P$\textup{)} has maximal depth, a kernel-$J$
fold or $iB$ depth $0$. Depth is \emph{order-respecting on the
nerve}: forgetting a factor can only increase depth,
$\operatorname{depth}(\partial_i\tau)\le\operatorname{depth}(\tau)$, since
$\overline{\partial_i\tau}$ receives a structure map from
$\overline\tau$, whence
$\Phi_{\overline\tau}\subseteq\Phi_{\overline{\partial_i\tau}}$ by
Lemma~\ref{lem:nervemaps}.
\end{lemma}

\begin{proof}
The map $\overline\tau\to\overline{\partial_i\tau}$ forgetting the
$i$-th factor is a structure map over $X$; a structure map
$W_1\to W_2$ over $X$ forces $\Phi_{W_1}\subseteq\Phi_{W_2}$
(Lemma~\ref{lem:nervemaps}). Both claims follow from
monotonicity of the kernel face along structure maps and the
definition of depth as poset-length of the kernel.
\end{proof}

\begin{remark}[Homotopy transfer along refinements]\label{rem:transfer}
If $R'$ refines $R$ with a choice of refinement map
$\sigma:R'\to R$ of families (each $R'$-piece factoring through
its $\sigma$-image), pullback of cochains along $\sigma$ is a
cochain map on nerves; two choices of $\sigma$ induce cochain
homotopic maps \textup{(}the standard comparison for \v{C}ech
refinements; the homotopy inserts at position $i$ the tuple
$(\sigma U_0,\dots,\sigma U_{i-1},U_i,\dots,U_q)$\textup{)}.
These comparison maps assemble the colimit of
Remark~\ref{rem:verdier}; in particular, a class vanishes in
sheaf cohomology iff it pulls to zero along \emph{some}
refinement, and any cocycle-level identity established after a
pullback persists in the colimit.
\end{remark}

\begin{definition}\label{def:square}
A cover of $X:=i(b)$ is \emph{standard} if all its piece maps are
surjective \textup{(}always achievable,
Lemma~\ref{lem:surj}\textup{)}, it is in the normal
form of Theorem~\ref{thm:normalform}, and moreover it contains the
pieces $iB$, $i(P/K)$ for all faces $K\supseteq J$, the dead
piece $D$, and two kernel-$J$ folds on distinct rays whenever
$\dim\bar P\ge2$ \textup{(}two nonproportional strictly positive
rational forms on the pointed cone $\Cone(\bar P)$; when
$\dim\bar P\le1$ no second fold is demanded, cf.\ the rank-one
comment in Step~2 of Proposition~\ref{prop:vanish}\textup{)}.
A standard cover is \emph{rigid} if no two distinct pieces admit
a structure map between them, except the quotient maps forced by
inclusion of faces. Every cover admits a standard rigid
refinement: first refine to normal form
(Theorem~\ref{thm:normalform}), then intersect objectwise with
the probe family $\{\mathfrak f\}\cup\{i(P/K)\}\cup\{iB\}\cup
\{D\}$ of Corollary~\ref{cor:probecovers}
\textup{(}possibly doubled to two rays; the intersection is a
cover, and all named pieces occur as their own
intersections\textup{)}; then delete, inductively, any piece with
a nonidentity structure map to another surviving piece---a point
served by the deleted piece is served via the composition, whose
birhizomicity at the point follows from the deleted map's
birhizomicity at the target factorization.
\textup{(}``Refined normal cover'' below means: standard rigid
cover.\textup{)}
\end{definition}

\begin{lemma}[Kernels in the probe nerve]\label{lem:nerve}
In the nerve of $\square_b$:
\textup{(}i\textup{)} $\Phi_{\,\mathfrak f\otimes P/K}=K$ for every
face $K\supseteq J$; in particular the overlap $\mathfrak
f\otimes iB$ has $B$-level $b$, and $\mathfrak f\otimes
P/K$ is dead for $K\supsetneq J$;
\textup{(}ii\textup{)} $\Phi_{P/K_1\otimes P/K_2}=K_1\vee K_2$;
\textup{(}iii\textup{)} consequently every nerve piece's $\mathcal
G$-value is $F$ of a single face quotient, with canonical
quotient maps between them.
\end{lemma}

\begin{proof}
(ii): quotients compose; by Proposition~\ref{prop:bp}(c) the joint
face has group $K_1^{\gp}+K_2^{\gp}$ and
$(K_1\vee K_2)^{\gp}\cap P=K_1\vee K_2$ refines the kernel: the
pushout of $P/K_1\leftarrow\bar P\to P/K_2$ is a quotient of both:
kernel exact face $K_1\vee K_2$ \checkmark.
(i): Lemma~\ref{lem:pushout} with $U=(P\to P/K)$, whence
$Z_U=P+K^{\gp}$, $L_U=K^{\gp}$: the jump condition reads: some
$t\ge1$, $a\in P^{\gp}$ with $z:=a-tx\in P+K^{\gp}$ and
$\chi(a)\le0$. Write everything downstairs:
$\bar z=\bar a-t\bar x\in\bar P+(K/J)^{\gp}$:
$\bar\chi(\bar a)-t\bar\chi(\bar x)=\bar\chi(\bar z)$; with
$\bar\chi$ tight: $\bar\chi(\bar a)\le0$; decompose
$\bar a=\bar\omega+\bar\kappa$ with $\bar\omega\in\bar P$ and
$\bar\kappa\in (K/J)^{\gp}$ \textup{(}possible since
$\bar a\in t\bar x+\bar P+(K/J)^{\gp}\subseteq
\bar P+(K/J)^{\gp}$\textup{)}. Then
$0\ge\bar\chi(\bar a)=\bar\chi(\bar\omega)+\bar\chi(\bar\kappa)$
with $\bar\chi(\bar\omega)\ge0$. Case A: $K=J$: $\bar\kappa=0$:
gives $\bar\chi(\bar\omega)=0$: $\bar\omega=0$ (tight): then
$\bar a=0$: jump condition was $-t\bar x\in\bar P$: sharpness:
$\bar x=0$: $x\in J$: so $\Phi=J$ \checkmark. Case B: $K\supsetneq J$:
$(K/J)^{\gp}\ne0$; since $\bar\chi$ is a group homomorphism,
$\bar\chi((K/J)^{\gp})=d\Z$ for some $d\ge1$ (it is nonzero:
otherwise $K^{\gp}\subseteq\ker\bar\chi^{\gp}$, hence
$K\cap\bar P=K\subseteq\ker\bar\chi=0$: contradiction with
$K\ne J$). Pick $\bar\kappa_N\in(K/J)^{\gp}$ with
$\bar\chi(\bar\kappa_N)=-Nd$. Then for ANY $\bar x\in\bar P$, with
$a$ defined by $\bar a:=t\bar x+\bar\kappa_N$ ($N$ large so
$\bar\chi(\bar a)\le0$), we get $\bar a-t\bar x=\bar\kappa_N\in
(K/J)^{\gp}\subseteq\bar Z$ \checkmark and $\bar\chi(a)\le0$ \checkmark: every $x$
dies: the piece is dead \checkmark. (For the second clause of (i):
$\mathfrak f\otimes iB$ is Case A.)
\end{proof}

\begin{lemma}[Kernels of general nerve pieces]\label{lem:mixed}
\textup{(}a\textup{)} The overlap of two kernel-$J$ folds has
kernel $J$ if their functionals are proportional
\textup{(}same ray\textup{)} and is dead otherwise.
\textup{(}b\textup{)} Longer overlaps are governed by the two
exact claims:
\begin{enumerate}[label=\textup{(}C\arabic*\textup{)}]
\item If $\Phi_U=J$ and $Z_U=P+J^{\gp}$ \textup{(}as is the case
for all pieces with surjective structure map,
Lemma~\ref{lem:Zpointed}\textup{)}, and $\varphi$ is tight on
$\bar P=P/J$, then $\Phi_{U\otimes\mathfrak f_\varphi}=J$.
\item If $\Phi_U=K$ and $Z_U=P+K^{\gp}$, and $\varphi$ vanishes
exactly on the face $K$ \textup{(}tight on $P/K$\textup{)},
then $\Phi_{U\otimes\mathfrak f_\varphi}=K$ \textup{(}applied
in the iteration with $K$ replaced by the finer
join\textup{)}.
\end{enumerate}
Consequently: if two positions are kernel-$J$ folds on distinct
rays, the piece is dead; otherwise its kernel equals the kernel
of the sub-overlap of its non-fold positions, computed
inductively by \textup{(}C2\textup{)}.
\end{lemma}

\begin{proof}
(a) Example~\ref{ex:ff}.
(C1) By Lemma~\ref{lem:pushout} a jump requires
$(t,a)$ with $a-tx\in Z_U=P+J^{\gp}$, i.e.\
$\bar a\in t\bar x+\bar P\subseteq\bar P$, and
$\varphi^{\gp}(a)\le0$. Since $\varphi$ is tight on $\bar P$,
$\bar a\in\bar P$ with $\varphi(\bar a)\le0$ forces
$\bar a=0$; then $-t\bar x\in\bar P$, and sharpness gives
$\bar x=0$: $x\in J$: no jump occurs, so the kernel stays $J$.
(C2) Identical after replacing $J$ by $K$: with
$\ker\varphi=K$, condition (B) gives
$\varphi^{\gp}(a)\le0$; write
everything modulo $J$, with $\bar\varphi:\bar P\to\mathbb N$
vanishing exactly on the face $K/J$: (A), i.e.\
$a-tx\in P+K^{\gp}$, gives
$\bar a\in t\bar x+\bar P+(K/J)^{\gp}$ after reduction;
absorbing the $(K/J)^{\gp}$-part into $\bar a$ (allowed:
$\bar\varphi$ vanishes there), (A) reads
$\overline{a-tx}\in\bar P$, i.e.\ $\bar a\in t\bar x+\bar P$;
(B) and the face property of $K/J$: if $\bar\varphi(\bar a)=0$
then $\bar a\in K/J$ and $t\bar x=\bar a-\bar z$ with
$\bar z\in\bar P$ is an equality inside $\bar P$ with sum in the
face $K/J$, hence $t\bar x\in K/J$; saturation gives
$\bar x\in K/J$, i.e.\ $x\in K$. If
$\bar\varphi(\bar a)<0$: $\bar\varphi(t\bar x)=
\bar\varphi(\bar a)-\bar\varphi(\bar z)\le\bar\varphi(\bar a)<0$
while $\bar\varphi(t\bar x)\ge0$: impossible. Conversely every
$x\in K$ dies via $t=1$, $\bar a=\bar x$ exactly as in the proof
of Theorem~\ref{thm:descent}, Claim B.
The final stronger clause is the iteration of (C1) (no fold
added) together with (C2) applied at the join kernel.
\end{proof}

\begin{lemma}[Value group on a face quotient]\label{lem:dlemma}
Let $K\supsetneq J$ be a face strictly larger than $J$, and let
$\chi:\bar P=P/J\to\mathbb N$ be tight. Then
$\bar\chi^{\gp}((K/J)^{\gp})=d\Z$ for some integer $d\ge1$
\textup{(}in particular it is nonzero\textup{)}.
\end{lemma}

\begin{proof}
Were it zero, $(K/J)^{\gp}\subseteq\ker\bar\chi^{\gp}$; applying
Proposition~\ref{prop:bp}(c) gives
$K/J=(K/J)^{\gp}\cap\bar P\subseteq\ker\bar\chi=\{0\}$ by
tightness, hence $K=J$: contradiction.
\end{proof}

\begin{lemma}[Constant strands]\label{lem:const}
Let $\mathscr W$ take value $A$ on every nerve piece of a cover
family $\mathcal I$, with every face map $\mathrm{id}_A$, and
assume the family contains two pieces $R_{\min}\subseteq
R_{\max}$ with $R_{\min}\subseteq U\subseteq R_{\max}$ for all
$U\in\mathcal I$. Then the alternating \v{C}ech complex computes
$\HH^0=A$ and $\HH^{q\ge1}=0$.
\end{lemma}

\begin{proof}
This is the constant-complex Lemma~\ref{lem:constcx} applied to
the family $\mathcal I$: $R_{\min}$ is its minimum piece, every
piece receives a structure map from it, and all face maps are
$\mathrm{id}_A$; the conclusion is
$\HH^0=A$, $\HH^{q\ge1}=0$ (the $\HH^0$-value is read off at any
piece, e.g.\ at $R_{\max}$).
\end{proof}


\begin{lemma}[The dead level]\label{lem:dead}
For the dead object $D=(P\to\{0\})$: $\mathcal G$ is constant equal
to $F(b_P)$ on its slice, every cover's pieces all lie over $D$, and
$H^0(D,\mathcal G)=F(b_P)$, $H^{q\ge1}(D,\mathcal G)=0$ by
Lemma~\textup{\ref{lem:const}} applied within each cover
\textup{(}equivalently: sheaf cohomology on a slice category with an
initial object, on a constant sheaf, vanishes in positive
degrees\textup{)}.
\end{lemma}

\begin{proof}
Every piece over $D$ has kernel $\supseteq\Phi_D=P$: all dead; the
slice $A_P/D$ has $D$ initial; compatible systems are pinned by
identity transitions to constants; the \v{C}ech complexes are the
constant-strand ones of Lemma~\ref{lem:const}.
\end{proof}

\begin{lemma}[Filter lemma]\label{lem:filter}
Let $\Xi$ be a finite poset of levels with a unique largest
element $\top$, and let $R$ be a cover of $X$ with a presheaf
$\mathscr W$ whose pieces carry order-respecting
$\Xi$-labels---every structure map
$\overline\tau\to\overline{\tau'}$ has
$\operatorname{lev}(\overline\tau)\XiLe
\operatorname{lev}(\overline{\tau'})$, so by
Lemma~\ref{lem:depth} each face of a nerve piece has label
$\XiLe$ the piece's label. Suppose:
\textup{(}F1\textup{)} the top level $\top$ is a
differential-closed constant strand with a minimum piece
\textup{(}Lemma~\ref{lem:constcx}\textup{)}: $\mathscr W$ is
constant equal to a single $A$ on $\top$-nerve pieces, all
structure maps among them inducing $\mathrm{id}_A$, and a fixed
piece $H_\top$ maps to every $\top$-piece;
\textup{(}F2\textup{)} every nerve piece labelled $\top$ arises
by adjoining to a lower-labelled tuple $\tau$ a piece $U$ whose
structure map factors through $H_\top$, with
$\overline{(U,\tau)}$ again $\top$-labelled; the piece $U$ is a
fold whose structure map's image lies in the kernel of the
transport $\mathscr W(\overline\tau\cup U)\to\mathscr
W(\overline\tau)$ \textup{(}the \emph{dead-adjunction
hypothesis}\textup{)}.
If $z$ is a degree-$q\ge1$ cocycle whose value at every
$\top$-labelled nerve piece vanishes, then $z$ is cohomologous,
within the same cover, to a cocycle supported on labels
$\XiLt\top$ \textup{(}the next-level values being the signed
transports of Lemma~\ref{lem:subquot}\textup{)}; iterating down
the finite poset, $z$ pulls \textup{(}after no refinement: at
the level of cocycles\textup{)} to a cocycle supported at the
initial level.
\end{lemma}

\begin{proof}
Define, for a lower tuple $\tau$ \textup{(}all labels
$\XiLt\top$\textup{)}, a cochain $h$ on $\top$-pieces by
$(he)_{(U,\tau)}:=e_{\tau}$ when (F2) applies
\textup{(}transported: $e_\tau$ lies at the lower level; the
value is its image at $\overline{(U,\tau)}$, well-defined since
the map $\overline{(U,\tau)}\to\overline\tau$ has constant
transport by (F2)'s factoring through $H_\top$
\textup{)}. Then the cocycle equation
$(\partial z)_{(U,\tau)}=0$ reads, isolating the face omitting
$U$,
\[
z_\tau\ \text{\textup{(}transported\textup{)}}\ =\
\sum_{i}(-1)^{i+1}\,\mathscr W(d_i)\bigl(z_{(U,\partial_i\tau)}\bigr)
\quad\text{on }\top\text{-pieces},
\]
whose right side is $(\partial h z)_{(U,\tau)}$: precisely the
constant-complex contraction of Lemma~\ref{lem:constcx} run
against the anchor $H_\top$. Subtracting $\partial(hz)$---a
coboundary supported on $\top$ by (F1)---removes the
$\top$-component, and by Lemma~\ref{lem:subquot} the cocycle is
now supported one level lower, the next-level values shifted by
the signed transport. Finiteness of $\Xi$ terminates the
descent at the initial level.
\end{proof}

\begin{lemma}[Sieve: fold terraces]\label{lem:sieve}
Let $R$ be a standard rigid cover of $i(b)$ and $z$ a $\mathcal
G$-\v{C}ech cocycle of degree $q\ge1$ on $R$. Then $z$ is
cohomologous within $R$ \textup{(}no refinement\textup{)} to a
cocycle vanishing at every nerve piece containing a kernel-$J$
fold \textup{(}the dead sector is left arbitrary\textup{)}.
\end{lemma}

\begin{proof}
A kernel-$J$ fold is a tight functional
$\varphi:\bar P\to\mathbb N$; up to rational multiple it is a
point of the pointed, strictly convex rational polyhedral cone
$\Cone(\bar P)^{\vee}$, so there are finitely many fold
\emph{rays} in $R$: enumerate them $\mathfrak f_1,\dots,\mathfrak
f_m$. For a ray index $r$ let the \emph{$r$-fold strand} consist
of the tuples containing $\mathfrak f_r$ but no fold of a
different ray. On such a tuple $\tau=(\mathfrak f_r,\sigma)$
every position of $\sigma$ is $iB$ or a fold on ray $r$, hence
has kernel $\subseteq J$; by Lemma~\ref{lem:nerve}(i) Case A and
Lemma~\ref{lem:mixed}(C1) the overlap
$\overline\tau=\mathfrak f_r\otimes\overline\sigma$ still has
kernel $J$, so all structure maps among these nerve pieces induce
the identity on the $\mathcal G$-value $F(b)$. Define a
$(q-1)$-cochain $w_r$ on the strand by
\[
(w_r)_\sigma\;:=\;z_{(\mathfrak f_r,\sigma)}\qquad
(\text{taken }=0\text{ if }(\mathfrak f_r,\sigma)\text{ has a
repetition}),
\]
which is well-defined because adjoining $\mathfrak f_r$ keeps the
kernel $J$ \textup{(}above\textup{)}. Then for a strand cocycle
$z$, $(\partial w_r)_\tau=z_\tau-(w_r\partial)_\tau$: this is the
constant-strand identity, since with all transports
$\mathrm{id}_{F(b)}$
\[
(\partial w_r)_\tau=\sum_i(-1)^i z_{(\mathfrak
f_r,\partial_i\tau)},\qquad
(w_r\partial)_\tau=z_\tau-\sum_i(-1)^i z_{(\mathfrak
f_r,\partial_i\tau)},
\]
so $(\partial w_r+w_r\partial)=z$ on the strand: the strand is
acyclic in degrees $\ge1$ and the $r$-component of $z$ is
$\partial w_r$. The strands for distinct $r$ are disjoint
\textup{(}a tuple with two distinct-ray folds lies in
none\textup{)} and their union is differential-closed, so by
repeated application of the exact-subquotient
Lemma~\ref{lem:subquot} the cocycle $z-\partial(\sum_r w_r)$ is
supported on tuples containing \emph{no} kernel-$J$ fold, as
claimed. The dead tuples are never touched.
\end{proof}
\begin{lemma}[Common refinement of covers]\label{lem:sandwich}
Let $R=\{U_\alpha\}$ and $S=\{V_\beta\}$ be covers of $i(b)$ in
the sliced birhizomic site. Then
$R^\flat:=\{U_\alpha\otimes_{i(b)}V_\beta\}_{\alpha,\beta}$ is a
cover of $i(b)$ refining both $R$ and $S$
\textup{(}via the structure maps to $U_\alpha$ and to
$V_\beta$\textup{)}; in particular the \v{C}ech colimit of
Remark~\ref{rem:verdier} is directed.
\end{lemma}

\begin{proof}
The birhizomic topology is a pretopology: its coverings are
stable under pullback and composition
\textup{(}Remark~\ref{rem:cover-points}; the pushouts exist by
Construction~\ref{constr:pushout}\textup{)}. For each fixed
$U_\alpha$ the family $\{U_\alpha\otimes_{i(b)}V_\beta\}_\beta$
is the pullback of the cover $S$ along $U_\alpha\to i(b)$, hence
a cover of $U_\alpha$; composing these over the cover
$R=\{U_\alpha\}$ of $i(b)$, the family
$\{U_\alpha\otimes_{i(b)}V_\beta\}_{\alpha,\beta}$ covers
$i(b)$. It refines $R$ and $S$ via the two structure maps
$U_\alpha\otimes V_\beta\to U_\alpha$ and
$U_\alpha\otimes V_\beta\to V_\beta$ given by the pushout
universal property.
\end{proof}


\begin{proposition}\label{prop:vanish}
\textup{(}a\textup{)} $H^0(i(b),\mathcal G)=F(b)$.
\textup{(}b\textup{)} For $q\ge1$, every class
$[z]\in\HH^{q}(R;\mathcal G)$ on a standard rigid cover is, after a
common refinement, cohomologous to a cocycle that \textup{(}i\textup{)}
has coboundary values with coefficients in $F(b_P)$ at every nerve
piece mapping to $D$, and \textup{(}ii\textup{)} vanishes at every
nerve piece containing a kernel-$J$ fold. The final vanishing
$H^{q>0}(i(b),\mathcal G)=0$ is \emph{not} proved for arbitrary $F$
\textup{(}see the Closing remark\textup{)}.
\end{proposition}

\begin{proof}
By the comparison theorem (Remark~\ref{rem:verdier}),
$H^q(i(b),\mathcal G)=\colim_R\HH^q(R;\mathcal G)$ over the nerves
of Definition~\ref{def:nerve}, the colimit taken over the standard
rigid covers (cofinal by Definition~\ref{def:square}), which form a
directed system by the common-refinement
Lemma~\ref{lem:sandwich}. We prove
$\HH^0(i(b),\mathcal G)=F(b)$---true already on the nose:
$\mathcal G$ is a sheaf (Theorem~\ref{thm:descent}) and
$\mathcal G(i(b))=F(b)$ (Proposition~\ref{prop:adj})---and that
every class $[z]\in\HH^{q\ge1}(R;\mathcal G)$ pulls to zero on a
further refinement. We work one degree $q\ge1$ at a time,
associating to a $q$-cocycle its $(q-1)$-\emph{faces} (the values
$z_{(U_0,U_1)}$), and use the sheaf axiom at each step to descend.
For $\bar P=\{0\}$ the slice is the dead level and
Lemma~\ref{lem:dead} gives the vanishing; for $\dim\bar P\le1$ the
whole nerve of a standard rigid cover is the constant strand with
minimum $iB$, so Lemma~\ref{lem:const}
($R_{\min}=R_{\max}=iB$) gives $\HH^{q\ge1}=0$ directly (these are
the base cases). Assume now $\dim\bar P\ge2$.

\emph{Step 1 (dead-sector hygiene).} Let $R$ be a standard rigid
cover. The sifted family $S:=\{U\otimes_{i(b)}D:U\in R\}$ is the
pullback of $R$ along $D\to i(b)$, hence a cover of $D$
(Lemma~\ref{lem:sandwich}); pass to the common refinement of $R$
and $S$ (Lemma~\ref{lem:sandwich} again) and keep the notation
$R$, $z$. Every nerve piece of $S$ is dead, and on the slice over
$D$ the sheaf $\mathcal G$ is the constant sheaf $F(b_P)$
(Lemma~\ref{lem:dead}); by the constant-complex
Lemma~\ref{lem:constcx} with anchor $H=D$, the restricted cocycle
is a coboundary: $z|_S=\partial w$ for a $(q-1)$-cochain $w$ on
$S$. Pulling $w$ back to $R$ along the refinement map, we may
subtract $\partial w$ from $z$; hence after this refinement
\[
\text{the value of }z\text{ at every nerve piece that maps to }D
\text{ is the value of a coboundary with coefficients in }F(b_P).
\tag{$\ast$}
\]
We assert nothing further about the dead sector; we do not pin it
to a single constant.

\emph{Step 2 (fold terraces).} Apply Lemma~\ref{lem:sieve} to
$z$: within the same cover it is cohomologous to a cocycle
vanishing at every nerve piece containing a kernel-$J$ fold;
replace $z$ by this representative. In particular
$z_{\mathfrak f}=0$ for every kernel-$J$ fold $\mathfrak f$ and
$z_\tau=0$ at every tuple containing one. \textup{(}The
identification $\bar\chi^{\gp}((K/J)^{\gp})=d\Z$, $d\ge1$, used
in the kernel computations is Lemma~\ref{lem:dlemma}.\textup{)}

\emph{Step 3 (where the argument stops).} After Step~2, $z$
vanishes on every tuple containing a kernel-$J$ fold; the surviving
tuples are built from $iB$ and the CFA pieces $i(P/K)$. The natural
candidate to conclude $[z]=0$ is the $iB$-contraction: adjoining
$iB$ is internal on a fold-free tuple, and
$(hz)_\tau:=z_{(iB,\tau)}$ \textup{(}$=0$ on repetition\textup{)} is
well-defined with $(\partial hz)_\tau=z_\tau-(h\partial z)_\tau$,
as in Lemma~\ref{lem:sieve}. However, to have
$z_\tau=(\partial hz)_\tau$ at a fold-free tuple one uses
$\partial z=0$ there, and the residual term of $\partial(hz)$ at the
dead adjacent tuples must then cancel the contraction: this requires
$\mathcal G$ to agree on the base face and its quotient faces, i.e.\
the transitions $\mathcal G(i(b))\to\mathcal G(i(P/K))$ to be
identities. For a general presheaf $F$ these are the quotient maps
$F(b)\to F(B(i(P/K)))$, which need not be identities, and the
discrepancy cannot be absorbed into Step~1's $w$
\textup{(}which lives only on the dead level\textup{)}. The
vanishing follows in those special cases where $\mathcal G$ agrees
on base and quotient faces, but not for arbitrary $F$; the document
therefore does \emph{not} assert $H^{q>0}(i(b),\mathcal G)=0$ in
general. The degree-$0$ part is Theorem~\ref{thm:descent} and
Proposition~\ref{prop:adj}.
\end{proof}

We close by recording exactly what has been proved and what remains
open, in accordance with the requirement never to present conjecture
as fact.

\begin{proposition}[The sheaf and the degree-zero unit]\label{thm:main-end}
For every finitely generated monoid $P$ and every presheaf $F$ of
abelian groups on $B_P^{\op}$:
\textup{(}i\textup{)} $\pi^{-1}F=G_F$ is a sheaf
\textup{(}Theorem~\ref{thm:descent}\textup{)}; \textup{(}ii\textup{)}
the unit $F\to\pi_*\pi^{-1}F$ is an isomorphism, i.e.\
$\mathrm R^{0}\pi_*(\pi^{-1}F)=F$.
\end{proposition}

\begin{proof}
(i) is Theorem~\ref{thm:descent}. (ii): with the trivial topology on
$B_P^{\op}$ sheaves equal presheaves, and
$H^0(i(b),\mathcal G)=\mathcal G(i(b))=F(b)$ by
Theorem~\ref{thm:descent} and Proposition~\ref{prop:adj}, so the
degree-zero unit $F\to\pi_*\pi^{-1}F$ is the identity sheaf map.
\end{proof}

\paragraph{Closing remark: higher direct images and the general
case.} Two points are deliberately left open, and this note does not
decide them. \textup{(}1\textup{)} \emph{Higher direct images.}
Proposition~\ref{prop:vanish} reduces $H^{q>0}(i(b),\mathcal G)$,
after the fold terraces are killed (Lemma~\ref{lem:sieve}) and the
dead sector is made coboundary-valued, to an $iB$-contraction on the
residuum. That contraction requires $\mathcal G$ to agree on the base
face and its quotient faces; for an \emph{arbitrary} presheaf $F$ the
structure map $\mathcal G(i(b))\to\mathcal G(i(P/K))$ is the quotient
map $F(b)\to F(B(i(P/K)))$, which need not be the identity, so the
residuum argument does not close $H^{q>0}(i(b),\mathcal G)=0$ for all
$F$, and the document does not claim it does. Consequently the full
quasi-isomorphism asked in the problem
\textup{(}$\mathrm R^{q>0}\pi_*(\pi^{-1}F)=0$ for $q>0$\textup{)} is
\emph{not proved here}. \textup{(}2\textup{)} \emph{Non-finitely
generated $P$.} As recorded in the Scope hypothesis
\textup{(}\S\ref{sec:intro}\textup{)} and Remark~\ref{rem:nonfg},
without finite generation of $P$ the map $P\to P/K$ need not be
finitely presented and the face poset can be infinite, so the
birhizomic-cover machinery and the finiteness used in the sieve
break; the general case is obstructed exactly there and is likewise
left open. No counterexample to either open point is known to this
note: as stated, the problem is \emph{open}, with the degree-zero
part and the sheaf property established above for finitely generated
$P$.

\end{document}
